\Needspace{4\baselineskip}
\begin{Example}{热方程}
\label{ex:heat-fundamental-solution}
对自然地作为初值问题提出的方程, 使用不同的基本解定义. 考虑方程
\begin{equation}\label{eq:first-order-time-evolution-fundamental-equation}
u_t(\bm{x},t)=L(D)u(\bm{x},t),\quad\bm{x}\in\mathbb{R}^m,\ t>0,
\end{equation}
其中 $L$ 为 $\mathbb{R}^m$ 上的常系数微分算子. 下文不将 $u$ 视为 $\mathbb{R}^m\times(0,\infty)$ 上的分布, 而将 $u$ 视为 $\mathbb{R}^m$ 上以 $t$ 为参数的分布. 若对每个 $\phi\in\mathcal{D}(\mathbb{R}^m)$, 积分
\begin{equation}\label{eq:time-parameter-distribution-test-pairing}
\int_{\mathbb{R}^m}u(\bm{x},t)\phi(\bm{x})\,\mathrm{d}\bm{x}
\end{equation}
关于 $t$ 连续, 则称 $u$ 连续依赖于 $t$; 若对每个 $\phi$, \eqref{eq:time-parameter-distribution-test-pairing} 都可微, 则称 $u$ 关于 $t$ 可微. 若 $u$ 关于 $t$ 可微, 则导数也是 $\mathbb{R}^m$ 上的分布; 这来自导数可表示为差商的极限以及分布空间的完备性.
\end{Example}

若 $u(\bm{x},t)$ 是 $\mathbb{R}^m$ 上连续依赖于 $t>0$ 的分布, 则也可以将 $u$ 视为 $\mathbb{R}^m\times(0,\infty)$ 上的分布. 这是因为每个试验函数 $\phi(\bm{x},t)\in\mathcal{D}(\mathbb{R}^m\times(0,\infty))$ 也可视为 $\mathbb{R}^m$ 上连续依赖于参数 $t$ 的试验函数. 由 Lemma~\ref{lem:distribution-test-sequence-pairing}, 这使
\begin{equation}\label{eq:time-dependent-test-distribution-pairing}
\int_{\mathbb{R}^m}u(\bm{x},t)\phi(\bm{x},t)\,\mathrm{d}\bm{x}
\end{equation}
成为 $t$ 的连续函数; 因而
\begin{equation}\label{eq:time-dependent-distribution-spacetime-pairing}
\int_0^\infty\int_{\mathbb{R}^m}u(\bm{x},t)\phi(\bm{x},t)\,\mathrm{d}\bm{x}\,\mathrm{d}t
\end{equation}
存在. 所以 $u$ 在 $\mathcal{D}(\mathbb{R}^m\times(0,\infty))$ 上定义一个线性泛函; 例如, 将 \eqref{eq:time-dependent-distribution-spacetime-pairing} 中的外层积分表示为 Riemann 和的极限, 并利用分布空间的完备性, 即可推得该泛函的连续性.

现在已准备好定义基本解.

\begin{Definition}
\label{def:evolution-fundamental-solution}
若 $G:[0,\infty)\to\mathcal{D}^{\prime}(\mathbb{R}^m)$ 在 $[0,\infty)$ 上连续可微, 且 $G$ 满足 \eqref{eq:first-order-time-evolution-fundamental-equation} 及初始条件 $G(\bm{x},0)=\delta(\bm{x})$, 则称 $G$ 为 \eqref{eq:first-order-time-evolution-fundamental-equation} 的\textbf{基本解}.
\end{Definition}

在这一定义中, 将 \eqref{eq:first-order-time-evolution-fundamental-equation} 中的 $u_t$ 视为关于参数 $t$ 的微分. 尽管如此, 容易证明 $G$ 在 $\mathbb{R}^m\times(0,\infty)$ 上的分布意义下也满足 \eqref{eq:first-order-time-evolution-fundamental-equation}.

非齐次问题
\begin{equation}\label{eq:first-order-evolution-inhomogeneous-initial-problem}
u_t=L(D)u+f(\bm{x},t),\quad u(\bm{x},0)=u_0(\bm{x}),
\end{equation}
的一个解, 其中 $f$ 和 $u_0$ 具有紧支集, 且 $f$ 从 $[0,\infty)$ 到 $\mathcal{D}^{\prime}(\mathbb{R}^m)$ 连续, 现在可以表示如下:
\begin{equation}\label{eq:first-order-evolution-fundamental-solution-formula}
\begin{aligned}
u(\bm{x},t)&=\int_{\mathbb{R}^m}G(\bm{x}-\bm{y},t)u_0(\bm{y})\,\mathrm{d}\bm{y}\\
&\quad+\int_0^t\int_{\mathbb{R}^m}G(\bm{x}-\bm{y},t-s)f(\bm{y},s)\,\mathrm{d}\bm{y}\,\mathrm{d}s.
\end{aligned}
\end{equation}
读者应验证这确实是一个解 (参见 Problem~\ref{ex:inhomogeneous-fundamental-solution-formula}).

现在考虑一维空间中的热方程. Problem~\ref{ex:heat-step-cauchy-problem} 要求求解问题
\begin{equation}\label{eq:heat-equation-heaviside-initial-problem}
\begin{aligned}
u_t&=u_{xx},\quad x\in\mathbb{R},\ t>0,\\
u(x,0)&=H(x),\quad x\in\mathbb{R}.
\end{aligned}
\end{equation}
可通过设 $u(x,t)=\phi(x/\sqrt t)$ 求得解, 这将问题化为常微分方程. 结果为
\begin{equation}\label{eq:heat-heaviside-initial-solution}
u(x,t)=\dfrac1{2\sqrt\pi}\int_{-\infty}^{x/\sqrt t}\exp\left(-\dfrac{v^2}4\right)\,\mathrm{d}v,
\end{equation}
为得到基本解, 只需关于 $x$ 求导. 于是得到
\begin{equation}\label{eq:heat-fundamental-solution-one-dimension}
G(x,t)=\dfrac1{2\sqrt{\pi t}}\exp\left(-\dfrac{x^2}{4t}\right).
\end{equation}
多维热方程的基本解可通过直积得到:
\begin{equation}\label{eq:heat-fundamental-solution-multiple-dimensions}
G(\bm{x},t)=\left(\dfrac1{2\sqrt{\pi t}}\right)^m\exp\left(-\dfrac{|\bm{x}|^2}{4t}\right).
\end{equation}

